\section{Introduction}

Long-horizon forecasting requires a model to extract predictive structure while the usefulness of past observations changes with the forecast length.
Modern approaches address this problem through efficient attention, decomposition, temporal patches, inverted tokenization, and multi-scale representations \cite{informer,autoformer,dlinear,patchtst,timesnet,itransformer}.
Their relative performance is strongly protocol- and dataset-dependent: lookback length, normalization, optimization budget, temporal split, and evaluation metric can alter model rankings \cite{m4,m5,tsf-survey}.

Nonstationarity adds a second difficulty.
Changes in level, scale, trend, volatility, or local dynamics can make future windows differ from the observations used for training.
RevIN addresses window-level distribution shift through reversible normalization, while Non-stationary Transformers combine stationarization with a learned mechanism intended to restore useful nonstationary information \cite{revin,nonstationary-transformer}.
Decomposition models provide a complementary inductive bias by separating slowly varying and residual components before forecasting \cite{autoformer,fedformer,dlinear}.

Mamba offers input-selective state-space sequence modeling with linear scaling in sequence length \cite{mamba}.
Forecasting adaptations differ in how they tokenize time and variables, scan in one or two directions, and combine temporal resolutions \cite{smamba,bimamba,mambats,timemachine,tsmamba}.
Recent multi-scale and decomposition-enhanced Mamba studies make this design space increasingly crowded \cite{msmamba,karma,dmamba}.
An efficient or more elaborate sequence operator is therefore not, by itself, evidence of improved forecasting.

The requirement for a strong reference forecast is particularly important for persistent financial level series.
Forecasting competitions repeatedly show that simple methods can remain difficult to beat \cite{m4,m5}, while financial forecasting reviews emphasize sensitivity to target definition, regime changes, and evaluation design \cite{finance-dl-survey}.
A model comparison on price or exchange-rate levels should therefore include a last-observation forecast before attributing an observed difference to architectural sophistication.

This study reports the completed repository benchmark for four Mamba configurations: S-Mamba, a decomposition variant (NS-Mamba-D), a multi-scale decomposition variant (MS-Mamba-MD), and a patch-temporal multi-scale decomposition variant (PMS-Mamba-PD).
They are compared with ten tuned neural baselines and a deterministic last-value forecast on GLD and Exchange under single-target (S-mode) forecasting.
Both datasets use 12 forecast lengths from 1 to 360 observations.
Within each final neural sweep, candidate configurations are selected by validation performance and evaluated on the test period.
However, earlier GLD test results were consulted during architecture development; consequently, the present evaluation is descriptive rather than an untouched confirmatory test.

The contributions are:
\begin{itemize}
    \item a traceable two-series benchmark covering 14 trained models, a deterministic last-value reference, and 12 forecast lengths;
    \item an implementation-grounded comparison of inverted, decomposition, multi-scale, and patch-temporal Mamba configurations against diverse neural alternatives;
    \item a reproducible persistence calculation showing that the strongest learned model differs by series and that last-value is best on Exchange under the current metric; and
    \item an explicit account of test-period reuse, single-seed uncertainty, overlapping forecast windows, search-budget asymmetry, post hoc dataset exclusion, and univariate tokenization.
\end{itemize}
